\documentclass[10pt,a4paper]{amsart}
\usepackage[margin=19mm]{geometry}
\usepackage{amsmath,amssymb,mathtools}
\usepackage[scr=rsfs,cal=euler]{mathalfa}
\usepackage{xfrac}
\usepackage[unicode,hidelinks,bookmarks=false]{hyperref}
\newcommand{\ie}{i.e.\ }
\newcommand{\cf}{cf.\ }
\newcommand{\resp}{respectively\ }
\newcommand{\Subsection}{\subsection}
\providecommand{\nobreakdash}{\nobreak\hbox{-}}
\setlength{\emergencystretch}{2em}
\multlinegap0pt
\usepackage{amssymb}
\usepackage{mathtools}
\usepackage{dsfont}
\usepackage{xcolor}
\usepackage{enumerate}
\newtheorem{definition}{Definition}
\newtheorem{proposition}{Proposition}
\newcommand{\am}{\alpha_-}
\newcommand{\ap}{\alpha_+}
\newcommand{\bigO}{\mathcal{O}}
\renewcommand{\div}{\operatorname{div}}
\newcommand{\dt}{\partial_t}
\newcommand{\eps}{\varepsilon}
\renewcommand{\epsilon}{\varepsilon}
\newcommand{\gm}{\gamma_-}
\newcommand{\gp}{\gamma_{+}}
\newcommand{\rmi}{\rho_-}
\newcommand{\rp}{\rho_+}
\title{Long title}
\author{Cassandre Lebot}
\date{Source: arXiv:2512.16286v2 (CC BY 4.0)}
\begin{document}
\maketitle

\noindent\emph{Source and licence.} Long title. Version arXiv:2512.16286v2 (CC BY 4.0), distributed by arXiv under the Creative Commons Attribution 4.0 International licence, http://creativecommons.org/licenses/by/4.0/, as recorded in the arXiv licence metadata for this exact version and shown on the abstract page https://arxiv.org/abs/2512.16286v2. Full licence notice: \texttt{licenses/arxiv-2512.16286-CC-BY-4.0.txt}.\par\smallskip

\noindent\textit{Editorial scope.} Counted unit: the article's proposition proposition (source0.tex lines 1070--1086). The article's complete original proof of it is delivered with it (source0.tex lines 1089--1149, one \texttt{proof} environment of 61 lines). No mathematics is written, repaired or re-derived: the statement and the proof are the article's own lines, in the article's own notation, and no retained line is substituted or re-flowed.  Retained uncounted as the source-local context the proof needs: lines 254--259; lines 843--845; lines 1053--1066.  Nothing was omitted from any retained span, so the package carries no omission record.  No retained line contains a citation, so the wrapper carries no bibliography and no bibliography entry.  No retained line contains a figure environment, an image-inclusion command or drawing code, so no image is delivered and the package declares no asset dependency.  Editorial formatting only, in addition: an AMS \texttt{amsart} preamble that restates the article's own declarations and macros used by the retained lines, namely the environments \texttt{definition}, \texttt{proposition} on separate counters, together with the standard packages those lines need.  The retained environments print in this document as Definition 1, Proposition 1 in order of appearance, whereas the article prints the same items with the numbering of its own full document.  The complete native source of the article as distributed in the arXiv e-print for this exact version is bundled unchanged as \texttt{sources/191231/source0.tex}.
\begin{definition}\label{defwpd} We say that the initial data $(u_0^\eps, \rho_0^\eps)$ are well-prepared if they satisfy
\begin{align*}
    \div u_0^\eps = \bigO(\eps^2), \quad  \rho_0^\eps = \bar{\rho_0} + \bigO(\eps^2)
\end{align*}
with $\bar{\rho_0}$ a constant.
\end{definition}

\begin{align}
    p_i(\rho_i,S_i) = p_i(\rho_i^0,S_i^0) + \eps^2 \Big( \partial_\rho p_i(\rho_i^0,S_i^0) \rho_i^1 + \partial_S p_i(\rho_i^0,S_i^0) S_i^1 + \cdots \Big).\label{expp}
\end{align}

We consider the following adimensional system:
\begin{align}
    \left\{
        \begin{array}{l}\label{PDE2veleps}
            \ap^\eps + \am^\eps = 1,\\
            \dt (\ap^\eps \rp) + \div(\ap^\eps \rp^\eps u_+^\eps) = 0,\\
            \dt (\am^\eps \rmi^\eps) + \div(\ap^\eps \rmi^\eps u_-^\eps) = 0,\\
            \dt (\ap^\eps \rp^\eps u_+^\eps) + \div (\ap^\eps \rp^\eps u_+^\eps \otimes u_+^\eps) + \frac{1}{\eps^2} \ap^\eps \nabla p_+(\rp^\eps, S_+^\eps) + p_{\mathrm{int}} \nabla \ap^\eps = 0,\\
            \dt (\am^\eps \rmi^\eps u_-^\eps) + \div (\am^\eps \rmi^\eps u_-^\eps \otimes u_-^\eps) + \frac{1}{\eps^2} \am^\eps \nabla p_-(\rmi^\eps, S_-^\eps) + p_{\mathrm{int}} \nabla \am^\eps = 0,\\
            \dt S_\pm^\eps + u_\pm^\eps \cdot \nabla S_\pm^\eps = 0,\\
            p_+(\rp^\eps, S_+^\eps) = p_-(\rmi^\eps, S_-^\eps).
        \end{array}
    \right.
\end{align}

\begin{proposition}
We consider a smooth enough solution of system \eqref{PDE2veleps} associated with well-prepared initial data in the sense of Definition \eqref{defwpd}, in a domain that is either periodic or bounded with the Dirichlet boundary conditions. We assume that both phases share the same low-Mach-number scaling, i.e., a single Mach number $\epsilon$ characterizes the flow. Taking the limit as $\eps$ goes to zero in the system \eqref{PDE2veleps}, we get formally the following limit system:
\begin{equation*}
    \begin{cases}
        \rp^0 = C^{\frac{1}{\gp}}e^{-S_+^0/\gp},\\
        \rmi^0 = C^{\frac{1}{\gm}}e^{-S_-^0/\gm},\\
        \div (\ap^0 u_+^0 +\am^0 u_-^0) = 0,\\
        \dt (\ap^0 \rp^0 u_+^0) + \div (\ap^0 \rp^0 u_+^0 \otimes u_+^0) + \ap^0 \nabla \pi^0 + p_{\mathrm{int}}^0 \nabla \ap^0 = 0,\\
        \dt (\am^0 \rmi^0 u_-^0) + \div (\am^0 \rmi^0 u_-^0 \otimes u_-^0) + \am^0 \nabla \pi^0 + p_{\mathrm{int}}^0 \nabla \am^0 = 0,\\
        \dt S_\pm^0 + u_\pm^0 \cdot \nabla S_\pm^0=0,\\
        \dt \alpha_+^0 + u_+^0 \cdot \nabla \alpha_+^0 = 0,\\
        \ap^0 + \am^0 = 1,\\
        \dt \alpha_+^0 + \div (\alpha_+^0 u_+^0) = 0.
    \end{cases}
\end{equation*}
with $\pi^0$ the limit pressure (of both phases).
\end{proposition}

\begin{proof}[Formal proof]
This formal proof is very similar to the one in the non-isentropic two-phase one-velocity case and we use the same expansion for the pressure \eqref{expp}. We obtain at order 1:
\begin{align*}
    \ap^0 + \am^0 = 1,\\
    \dt S_\pm^0 + u_\pm^0 \cdot \nabla S_\pm^0 = 0.
\end{align*}
and the algebraic closure reads
\begin{align*}
    p_+ (\rp^0,S_+^0) = p_-(\rmi^0,S_-^0).
\end{align*}

At order $\eps^{-2}$ the momentum equations yield
\begin{align*}
    \left\{
        \begin{array}{l}
              \nabla p_+(\rho_+^0) = 0, \\
              \nabla p_-(\rho_-^0) = 0.
        \end{array}
    \right.
\end{align*}

Thus, we can set $C(t)$ a constant in space such that
\begin{align*}
     p_+ (\rp^0,S_+^0) = p_-(\rmi^0,S_-^0) = C(t),
\end{align*}

Adding the two mass equations using the entropy equations and the state laws, we now obtain
\begin{align}
    \div (\ap^0 u_+^0 + \am^0 u_-^0) + \Big( \frac{\ap^0}{\gp} + \frac{\am^0}{\gm} \Big) \frac{C'(t)}{C(t)}=0.\label{divuI2}
\end{align}

Now we integrate in space, and because of our assumption on $\Omega$, we get
\begin{align*}
    \frac{C'(t)}{C(t)}\int_\Omega \Big( \frac{\ap^0}{\gp} + \frac{\am^0}{\gm} \Big) d x   =0.
\end{align*}

Thus $C'(t) = 0$ and $C$ is a constant in time and space. Coming back to \eqref{divuI2} we obtain
\begin{align*}
    \div (\ap^0 u_+^0 + \am^0 u_-^0) = 0.
\end{align*}



The mass equations at order $1$ now read
\begin{align*}
    \dt \alpha_\pm^0 + \div (\alpha_\pm^0 u_\pm^0)=0.
\end{align*}

Setting $\pi_i^0 = p_i(\rho_i^0) \Big( \gamma_i \frac{\rho_i^1}{\rho_i^0} + S_i^1 \Big)$ for $i=+,-$, we get from the order $\eps^2$ of the algebraic closure:
\begin{align*}
    \partial_\rho p_+(\rp^0,S_+^0)\rp^1 + \partial_S p_+ S_+^1 &= (\partial_\rho p_- \rmi^1 + \partial_S p_- S_-^1),\\
    \pi_+^0 &= \pi_-^0 := \pi^0. 
\end{align*}

Finally, the order 1 momentum equations read
\begin{align*}
    \rp^0 \ap^0 ( \dt u_+^0 + u_+^0 \nabla \cdot u_+^0) + \ap^0 \nabla \pi^0 + p_{\mathrm{int}}^0 \nabla \ap^0 = 0,\\
    \rmi^0 \am^0 (\dt u_-^0 + u_-^0 \nabla \cdot u_-^0) + \am^0 \nabla \pi^0 + p_{\mathrm{int}}^0 \nabla \am^0 = 0.
\end{align*}

\end{proof}

\end{document}
